\section{Symplectic Integrators}

\subsection{Motivation}

Numerically integrating the equations of motion of an orbiting body is a non-trivial task. A naïve approach — such as the simple Euler method — accumulates energy errors over time, causing simulated orbits to either spiral inward or drift outward, even though the true Newtonian orbit is a closed, energy-conserving ellipse.

A better class of methods, known as \emph{symplectic integrators}, preserves the geometric structure of Hamiltonian mechanics. Crucially, they conserve a quantity that is very close to the true total energy over long integration times, preventing the secular drift that plagues non-symplectic methods. The \emph{Velocity Verlet} algorithm is the simplest and most widely used symplectic integrator, and is the method employed in this simulation.

\subsection{The Velocity Verlet Algorithm}

The algorithm advances the state of the system — position $\vec{r}$, velocity $\vec{v}$, and acceleration $\vec{a}$ — by a fixed timestep $dt$ in three sequential sub-steps.

\subsubsection{Step 1: Position Update}

The new position at timestep $n+1$ is computed using the current position, velocity, and acceleration:

\begin{equation}
\vec{r}_{n+1} = \vec{r}_n + \vec{v}_n\,dt + \frac{1}{2}\,\vec{a}_n\,dt^2
\end{equation}

This is equivalent to a second-order Taylor expansion of the position. The $\frac{1}{2}\vec{a}\,dt^2$ term distinguishes Verlet from the simple Euler method ($\vec{r}_{n+1} = \vec{r}_n + \vec{v}_n\,dt$) and provides second-order accuracy in the position update.

\subsubsection{Step 2: Acceleration Update}

With the new position $\vec{r}_{n+1}$ known, the updated gravitational acceleration is recalculated from the force law:

\begin{equation}
\vec{a}_{n+1} = -\frac{GM}{|\vec{r}_{n+1} - \vec{r}_\star|^3}\,(\vec{r}_{n+1} - \vec{r}_\star)
\end{equation}

This step uses the \emph{updated} planet position but the \emph{fixed} star position, correctly computing the new gravitational pull after the planet has moved.

\subsubsection{Step 3: Velocity Update}

Finally, the velocity is updated using the \emph{average} of the old and new accelerations:

\begin{equation}
\vec{v}_{n+1} = \vec{v}_n + \frac{1}{2}\,(\vec{a}_n + \vec{a}_{n+1})\,dt
\end{equation}

Using the mean acceleration over the interval, rather than just the initial acceleration (as Euler would), is what gives Velocity Verlet its superior energy conservation properties.

\subsection{Why Symplectic Integration Matters}

The defining property of a symplectic integrator is that it preserves the \emph{symplectic structure} of the Hamiltonian phase space, which in practice means the total mechanical energy of the simulated system does not drift monotonically over time. Instead, it oscillates around the true value by a small, bounded amount. For an orbital simulation run over thousands or millions of steps, this is essential: a non-symplectic method might produce a visually plausible orbit for a few hundred steps before the energy error accumulates enough to deform or destroy the trajectory.

\subsection{Implementation Summary}

The simulation runs the Velocity Verlet loop for 10,000 timesteps with $dt = 1$. At each step the following operations are performed in order:

\begin{enumerate}
    \item Compute new planet position from $(\vec{r}_n,\, \vec{v}_n,\, \vec{a}_n)$.
    \item Recompute the planet-star distance from the new position.
    \item Compute new acceleration from the updated distance and position.
    \item Compute new velocity from both old and new accelerations.
    \item Record the snapshot $(\vec{r}_{n+1},\, \vec{v}_{n+1},\, \vec{a}_{n+1})$ for CSV export.
\end{enumerate}

After all 10,000 steps, the trajectory data is written to \texttt{orbit\_data.csv} and visualised via the accompanying Python plotter, producing the orbital path shown in the output figure.